\section{Conditional SAGE and packed development comparison}\label{app:E-PE}
The conditional comparison tests a change to the neighbor values while retaining native SAGE aggregation. All 15 banks and 24 optimizer/body acquisitions in E closed before the complete readout; all 75 temperature fits closed. A separate packed comparison adds three banks, three optimizer bundles and 12 private narrower bodies. The joined raw analysis therefore covers 18 banks, 27 acquisition bundles and 36 body records, without adding calibration, model forwards or selected-state deserializations.

\subsection{Operation, controls and acquisition}
All arms use the complete WikiCS graph (11,701 nodes, 300 features and ten classes), 580 TRAIN nodes and 5,274 development nodes, with seeds 7301, 7403 and 7507. BASE/E4/N4/equipped M1/I4 use two native residual SAGE blocks, width 128, LayerNorm and dropout 0.2. Internal BatchEnsemble factors wrap every native dense map; shared native biases remain outside the private output scales. Native first-input factors start with Rademacher signs, later factors at one, and each route has a persistent dropout stream.

For each factual nonself edge $u\to v$ and route $m$, E4's two $256\to16$ affine builders consume $z_{uv}^{(m)}=[h_u^{(m)},h_v^{(m)}]$. Their shared matrices and biases have private input/output factor vectors. Reshaping each output to $4\times4$ gives
\begin{equation}\label{eq:E-paired-maps}
F_{a,uv}^{(m)}=I_4+\tfrac12\tanh A_a^{(m)}(z_{uv}^{(m)}),\quad a\in\{\mathrm{src},\mathrm{tgt}\},\qquad
T_{v\leftarrow u}^{(m)}=F_{\mathrm{tgt},uv}^{(m)\top}F_{\mathrm{src},uv}^{(m)}.
\end{equation}
The product acts on the $4\times32$ source-value view before scalar mean aggregation. Self values, support and native root/neighbor affine maps, feed-forward order, normalization, dropout and classifier are retained. Builder weights/biases start at zero, so transport starts at identity; private builder factors start at one and acquire gradients after the builder weights move. This initialization supplies no learning guarantee. N4 uses $[h_v^{(m)},\operatorname{SiLU}(h_v^{(m)})]$ with the same builders and matrix action, built once per receiver rather than once per edge.

Shared banks minimize mean own cross-entropy; no pooled-loss gradient or auxiliary loss is added. AdamW uses learning rate 0.001 and zero decay, at most 1,000 updates and 300 stale updates. The selector keeps the earliest strict full-development accuracy maximum of the probability mean and restores model, optimizer and route streams. Equipped M1 has its own full loss and selector. Equipped I4 comprises four separately initialized, optimized, selected and restored M1 bodies with the same endpoint operation, followed by probability pooling. It has no reused M1 fit or learned-state cloning. These are complete acquisition procedures, with shared-bank and independent-body optimization/selection histories.

PE4 is a declared native SAGE adaptation of Packed-Ensembles with four untied width-59 bodies and $\gamma=1$. Dense maps, native biases, LayerNorms and heads are private, while parameter-free neighbor means are packed over member channels. Each body has complete hidden states and gradients. Mean own CE and one coherent pooled selector follow the same F training contract. Width is chosen by the full BASE parameter cap before scoring: width 59 gives 215,508 parameters; width 60 would give 221,560 and exceed the 216,674 cap. E4 has 21,966 more parameters than this packed reference. This fixed narrower context leaves the full-width equipped I4 quality comparison in place. The earlier F outcomes and original benchmark are retained separately; no historical cross-runtime F score supplies these fresh anchors.

\subsection{Complete selected outcomes}

\begin{table}[t]
\centering\small\setlength{\tabcolsep}{4pt}
\caption{Complete joined raw WikiCS development outcomes. Accuracy and member means are percentages, averaged over three optimizer seeds. NLL and Brier are lower-is-better. Every bank uses the native checkpoint-selection population.}
\label{tab:E-PE-raw}
\begin{tabular}{lrrrrr}
\toprule
Arm & Accuracy & Member mean & NLL & Brier & Parameters\\
\midrule
BASE (fresh shared4) & 80.6282 & 78.3766 & 1.137994 & 0.315096 & 216,674\\
E4 (endpoints) & 81.1971 & 79.4495 & 0.761107 & 0.297086 & 237,474\\
N4 (receiver) & 81.0517 & 79.5174 & 0.737769 & 0.295872 & 237,474\\
Equipped M1 & 80.2174 & 80.2174 & 0.765341 & 0.314952 & 225,600\\
Equipped I4 & 80.8937 & 80.2127 & 0.688787 & 0.293094 & 902,400\\
PE4 ($\gamma=1$, width 59) & 79.7876 & 78.4936 & 0.695566 & 0.306853 & 215,508\\
\bottomrule\end{tabular}\end{table}

Calibration divides the development nodes into five label-free folds of sizes 1,055/1,055/1,055/1,055/1,054, after native checkpoint selection used all 5,274 nodes. For each arm/seed/fold, one CPU float64 temperature fit completes 500 updates, with no retries or final whole-development refit. All 75 fits and 37,500 scalar updates closed. Out-of-fold calibration consequently does not create independent whole-pipeline evaluation. Calibration preserves member argmaxes but can change the probability mean. No PE temperature or calibrated PE endpoint was fitted.

\begin{table}[t]
\centering\small\setlength{\tabcolsep}{5pt}
\caption{Complete E-family calibrated development outcomes. Every entry is a three-seed mean. PE is absent because the joined comparison added no calibration.}
\label{tab:E-calibrated}
\begin{tabular}{lrrr}
\toprule
Arm & Calibrated accuracy (\%) & NLL & Brier\\
\midrule
BASE (fresh shared4) & 80.7294 & 0.685842 & 0.292675\\
E4 (endpoints) & 81.1971 & 0.652050 & 0.284499\\
N4 (receiver) & 81.0201 & 0.649441 & 0.284886\\
Equipped M1 & 80.2174 & 0.678221 & 0.297737\\
Equipped I4 & 80.9506 & 0.636007 & 0.283364\\
\bottomrule\end{tabular}\end{table}

\begin{table}[t]
\centering\small\setlength{\tabcolsep}{5pt}
\caption{E4 minus each E-family reference. Raw accuracy differences and descriptive paired-seed 95\% intervals are in percentage points; calibrated NLL differences are signed, with positive values worse. All raw seed differences are positive.}
\label{tab:E-paired}
\begin{tabular}{lrrr}
\toprule
Reference & Raw gain & Descriptive 95\% interval & Calibrated NLL change\\
\midrule
BASE & 0.5688 & $[0.1146,1.0231]$ & $-0.033791$\\
Equipped I4 & 0.3034 & $[-0.3560,0.9628]$ & $+0.016044$\\
N4 & 0.1454 & $[0.0015,0.2893]$ & $+0.002609$\\
Equipped M1 & 0.9796 & $[0.2124,1.7469]$ & $-0.026171$\\
\bottomrule\end{tabular}\end{table}

\paragraph{Predeclared screens and uncertainty.}
The E raw rule requires mean gains of at least 0.2 points against BASE and equipped I4, all three differences nonnegative and at least two positive. Endpoint attribution against N4 has the same sign rule and a mean margin of 0.1 points. Both accuracy screens pass. The separate calibrated-NLL guard allows mean harm at most 0.02 and each-seed harm at most 0.05 against BASE/I4; it also passes. Against I4, the three calibrated harms are 0.010782/0.024975/0.012375, with mean 0.016044 and descriptive interval $[-0.003271,0.035359]$. All three are harm despite passing the practical allowance. Raw and calibrated NLL and Brier favor I4. Calibrated accuracy favors E4 at two seeds, with mean gain 0.2465 points and descriptive interval $[-0.3712,0.8642]$.

These intervals use three paired optimizer repeats on one encountered graph ($t$, two degrees of freedom). They do not account for preceding design selection, multiple comparisons or variation across graphs. The raw exact two-sided sign-flip value is 0.25 for each E-family contrast. Node readouts and calibration folds are not independent graph trials. The I4 accuracy interval includes zero; a practical NLL allowance does not establish statistical noninferiority. Member quality and cost were diagnostics rather than post hoc vetoes. No test labels enter this development analysis, and no unused confirmation or novelty verdict follows.

\begin{table}[t]
\centering\small\setlength{\tabcolsep}{4pt}
\caption{All joined packed-context raw contrasts, candidate minus PE4. Gains and descriptive paired 95\% intervals are in percentage points. Positive NLL/Brier changes are worse. These contexts do not replace the E-family's predeclared full-width quality references.}
\label{tab:PE-contrasts}
\begin{tabular}{lrrrr}
\toprule
Candidate & Accuracy gain & Descriptive 95\% interval & NLL change & Brier change\\
\midrule
BASE (fresh shared4) & 0.8406 & $[0.7318,0.9494]$ & $+0.442429$ & $+0.008244$\\
E4 (endpoints) & 1.4094 & $[0.9820,1.8368]$ & $+0.065541$ & $-0.009767$\\
N4 (receiver) & 1.2641 & $[0.7560,1.7721]$ & $+0.042203$ & $-0.010981$\\
Equipped M1 & 0.4298 & $[-0.0228,0.8824]$ & $+0.069775$ & $+0.008099$\\
Equipped I4 & 1.1061 & $[0.8643,1.3478]$ & $-0.006779$ & $-0.013758$\\
\bottomrule\end{tabular}\end{table}


The BASE--PE and E4--PE raw screens pass the declared 0.2-point mean/sign rule. PE improves NLL over BASE/E4/N4/M1 but has worse NLL than I4; its Brier ordering also differs. One width/gamma choice supplies no uniform quality or whole-family ranking.

\subsection{Prediction accounting}
Covered alternatives are readouts with at least one correct member. A pooling loss is a covered alternative lost in the served probability mean. The counts below sum three dependent readouts of the same 5,274 nodes (15,822 readouts) and compare separately fitted/selected procedures; they are descriptive rather than causal events.

\begin{table}[t]
\centering\small\setlength{\tabcolsep}{4pt}
\caption{Complete net error accounting. Each row is candidate minus reference. Coverage and loss changes retain their signs; an aggregation-only correction has no individually correct member. Repairs/harms compare the two pooled decisions directly.}
\label{tab:E-PE-errors}
\begin{tabular}{lrrrrr}
\toprule
Contrast & Coverage & Pool losses & Pool-only & Net correct & Repairs/harms\\
\midrule
E4--BASE & $-106$ & $-196$ & 0 & 90 & 512/422\\
E4--N4 & $-12$ & $-35$ & 0 & 23 & 213/190\\
E4--equipped I4 & 44 & $-4$ & 0 & 48 & 475/427\\
E4--equipped M1 & 970 & 815 & 0 & 155 & 595/440\\
\midrule
BASE--PE4 & $+295$ & $+161$ & $-1$ & 133 & 773/640\\
E4--PE4 & $+189$ & $-35$ & $-1$ & 223 & 784/561\\
N4--PE4 & $+201$ & $+0$ & $-1$ & 200 & 729/529\\
Equipped M1--PE4 & $-781$ & $-850$ & $-1$ & 68 & 553/485\\
Equipped I4--PE4 & $+145$ & $-31$ & $-1$ & 175 & 531/356\\
\bottomrule\end{tabular}\end{table}

E4 improves members over BASE while coverage falls and fewer alternatives are lost in pooling; N4 serves 67 of E4's 90 additional BASE-correct readouts. The endpoint increment varies across classes. These are separately fitted/selected procedures, without private-builder or optimization attribution. Complete discussion is in \path{code_evidence/native_E_PE/complete_contrasts_and_accounting.tex.txt}.

\subsection{Recorded resource scopes}
\begin{table}[t]
\centering\small\setlength{\tabcolsep}{5pt}
\caption{Complete three-bank resource scopes. Acquisition and selected serving are sums in seconds; allocated peak is decimal GB. I4 peak is the maximum for a separately acquired complete member body. Timers include different stopping histories and overlap; they are not isolated latency benchmarks.}
\label{tab:E-PE-costs}
\begin{tabular}{lrrr}
\toprule
Arm & Acquisition sum & Selected serving sum & CUDA allocated peak (GB)\\
\midrule
BASE (fresh shared4) & 234.671 & 0.161034 & 1.162\\
E4 (endpoints) & 1566.748 & 1.325601 & 16.120\\
N4 (receiver) & 636.351 & 0.311309 & 4.018\\
Equipped M1 & 372.650 & 0.351277 & 4.652\\
Equipped I4 & 1487.302 & 1.293800 & 4.652\\
PE4 ($\gamma=1$, width 59) & 134.589 & 0.126955 & 1.020\\
\bottomrule\end{tabular}\end{table}

Sharing saves stored parameters against equipped I4, while E4's observed four-route activation peak is much larger than that of one separately acquired I4 body. The E acquisition overlapped other work on the same GPU, and the PE acquisition was a separate run. Stopping histories, selector calls, restores and serving scopes are retained; sums are not concurrent wall time, and lifetime RSS peaks are not per-arm allocations. Preparation and discarded qualification also incur cost. No lower compute, isolated latency or peak-memory claim follows from the parameter counts.

All original E calibration rules remain unchanged. PE adds raw endpoints only. The operation combines established mechanisms; the comparison does not isolate private builder factors from private route states or capacity.

\subsection{Error strata and anchored head diagnostic}\label{app:E-error-head}
A fixed diagnostic reads the archived E/BASE/N4/I4 decisions at seeds 7301/7403/7507 without changing them. The 15,822 readouts repeat 5,274 development nodes. A strict common rival is some false class strictly above truth in every member's raw scores, rather than a shared wrong argmax. Exact softmax preserves each strict ordering; any nonnegative scalar probability average preserves their common ordering. This limits unchanged member mixing, without proving absent hidden information or locating a representation/classifier bottleneck.

E4 has 2,064 strict-rival readouts, versus BASE's 1,929 and I4's 2,139. Of E4's 2,975 pooled errors, 815 lose a correct member alternative and 96 have neither a correct member nor a strict common rival. Against BASE, E4 removes 316 old rival predicates, serves 79 old-rival repairs and creates 451; against I4 it removes 436, serves 122 and creates 361. Predicate removal alone is not served repair. E4--I4's net +48 consists of +63/+45 for factual classes 2/8 and -60 elsewhere; five classes gain and five lose. E4--N4's +23 consists of +27 at incoming nonself degree 5--12 and -4 elsewhere. These post hoc marginals justify no class/degree router, monotone degree effect or causal source-context claim. Full cells, including empty/small strata, remain in \path{evidence/closed_E_strata_COUNTS.json}.

The completed head diagnostic holds the original eval-mode width-128 representations and original heads fixed for affine residual conditions. Across 18 restored E/M1/I4 bodies, 27 member trajectories retain all 580 TRAIN and 5,274 SELECT rows. E compares head-only BatchEnsemble continuation with common and private zero-initialized unrestricted affine residuals. M1/I4 receive equally processed continuation/affine controls; I4 retains four own optimizers/selectors. Every condition adds a learned score delta to the original stored-logit anchor, and original flags remain authority. Fresh AdamW moments, rate 0.001, zero decay, mean own CE, a 1,000-update cap and 300 stale updates reproduce the strict full-SELECT probability-mean selector. Steps 1--1,000 are eligible; untouched step 0 remains a separate reported anchor. There is no backbone update, calibration, feature scaling or class-specific fitting.

\begin{table}[t]
\centering\small\setlength{\tabcolsep}{4pt}
\caption{Complete anchored fixed-feature head diagnostic. Values are equal means across three optimizer seeds; accuracy/member means are percentages and raw pooled NLL is lower-is-better. Untouched references remain visible. All outcomes use the already selected development population.}
\label{tab:E-head-absolute}
\begin{tabular}{lrrrr}
\toprule
Condition & Accuracy & Member mean & Raw NLL & Strict rivals\\
\midrule
E4 untouched & 81.197067 & 79.449501 & 0.761106833 & 2064\\
E4 BE continuation & 81.178106 & 79.479522 & 0.777546617 & 2059\\
E4 common affine & 81.159145 & 79.449501 & 0.774839515 & 2061\\
E4 private affine & 81.190747 & 79.593288 & 0.780286927 & 2080\\
M1 untouched & 80.217419 & 80.217419 & 0.765340731 & 3130\\
M1 BE continuation & 80.299583 & 80.299583 & 0.772817430 & 3117\\
M1 affine & 80.312223 & 80.312223 & 0.771618117 & 3115\\
I4 untouched & 80.893692 & 80.212679 & 0.688786634 & 2139\\
I4 BE continuation & 80.912653 & 80.263241 & 0.694397369 & 2143\\
I4 private affine & 80.906333 & 80.271141 & 0.693524059 & 2142\\
\bottomrule\end{tabular}\end{table}

All 21 banks and 39 head optimizer bundles close. The frozen private-specific screen requires positive whole-SELECT private-minus-common and private-minus-BE accuracy at every seed, no worse mean member competence/NLL and positive mean improvement over untouched E. Correct-count vectors are $+3,-1,+3$ versus common, $+3,-3,+2$ versus BE and $+2,-3,0$ versus untouched E. Every accuracy clause fails. Mean member accuracy improves by 0.143787/0.113766 points versus common/BE, but pooled NLL worsens by 0.005447412/0.002740309 and by 0.019180094 versus untouched E. Per-member SELECT NLL was not exported for E/I4 and is unevaluable; pooled NLL does not substitute for it.

Private E has 50 own-anchor repairs and 51 harms. Coverage adds 45/removes 56 and pooling losses fall by 10, giving net -1 correct. On all 2,064 original E rival errors it obtains 40 correct-member alternatives and removes 46 predicates, yet loses all 40 in pooling and serves zero repairs. It creates 62 whole-population rival predicates, ending at 2,080. Common/BE also serve zero old-E-rival repairs. Private E remains above untouched M1/I4 by 0.973328/0.297055 points, but these advantages predate the head change; against affine-processed M1/I4 it gains 0.878524/0.284414 points, with I4 correct-count signs $+24,+27,-6$. Its NLL is worse than every original/processed reference. Private-minus-untouched-E class nets are $3,0,-21,8,-9,4,1,12,3,-2$, summing to -1; no favorable class rescues the failed screen.

The recorded complete scope is 74.54 seconds, process-lifetime peak RSS 943,833,088 bytes and CUDA allocation 1,569,668,608 bytes, with feature extraction/head fitting overlapping other work. These are not isolated speed or per-arm peaks. The earlier negative regularized Polynormer head-refit scope remains closed. This result closes the exact anchored affine recipe, without ruling out nonlinear information, stronger ordinary/end-to-end training or global capacity. Full protocol, absolute per-seed records, contrasts, factual classes, optimizer accounting and cost/replay receipts are in \path{evidence/closed_E_head_PLAN.json}, \path{evidence/closed_E_head_COMPLETE.json} and \path{evidence/closed_E_head_SUMMARY.json}. Numerical replay is an engineering record rather than a utility criterion; no unused whole-pipeline confirmation follows.
